\documentclass[12pt]{article}
\usepackage[T2A]{fontenc}
\usepackage[utf8]{inputenc}
\usepackage[russian]{babel}
\usepackage{amsmath,amssymb,amsthm}
\usepackage{geometry}
\geometry{margin=2.5cm}
\usepackage{indentfirst}

\title{Спектральная геометрия фрактальных графов\\ и аномальная диффузия на них}
\author{Е.~Р.~Громова}
\date{2024}

\begin{document}
\maketitle

\section*{Аннотация}
Вводится класс фрактальных графов, возникающих как пределы итеративных функциональных систем с переменным ветвлением. Вычислена спектральная размерность и показано, что случайное блуждание на таких графах демонстрирует субдиффузию с показателем, зависящим от мультифрактального спектра.

\vspace{0.5cm}
\noindent\textbf{Ключевые слова:} фрактальные графы; спектральная размерность; аномальная диффузия; мультифрактальный анализ; случайные блуждания.

\vspace{0.3cm}
\noindent\textbf{Раздел журнала:} Математика. \quad \textbf{Файл:} 1 из 2.

\section{Введение}
Фрактальные графы служат дискретными моделями неоднородных сред. Цель работы — установить связь между спектральной геометрией и характером диффузии.

\section{Основной результат}
Пусть $d_s$ — спектральная размерность графа. Тогда среднеквадратичное смещение блуждания удовлетворяет
\begin{equation}
\langle r^2(t) \rangle \sim t^{2/d_w}, \qquad d_w = \frac{2 d_s}{d_s - 1},
\end{equation}
где $d_w$ — показатель аномальной диффузии.

\section{Заключение}
Полученные результаты применены к моделированию транспорта в пористых средах.

\end{document}